\begin{figure}[H]\centering\resizebox{0.98\linewidth}{!}{%
\begin{tikzpicture}[
 b/.style={draw,rounded corners=1.2pt,align=center,text width=3.75cm,minimum height=0.92cm,fill=blue!6},
 d/.style={draw,diamond,aspect=2.25,align=center,inner sep=2pt,fill=orange!7},
 a/.style={-Stealth,thick},font=\small]
\node[b,text width=4.0cm] (access) at (-3.0,4.8) {EXU 发起数据访问};
\node[d] (kind) at (3.0,4.8) {按 tag 结果与访问类型分派};
\draw[a] (access)--(kind);

% 四条等宽泳道使四种机制可独立顺读。
\coordinate (busc) at (3.0,3.55);
\draw[thick] (kind)--(busc);
\draw[thick] (-7.2,3.55)--(8.4,3.55);

\node[b,fill=green!7] (normal) at (-7.2,2.0) {普通 load 命中\\主 BRAM 在 C1 同步读出};
\node[b,fill=green!7] (lsu) at (-7.2,0.0) {数据、地址与 bank 对齐\\LSU 扩展并前递};
\draw[a] (-7.2,3.55)--(normal); \draw[a] (normal)--(lsu);

\node[b,fill=cyan!8] (mirror) at (-2.0,2.0) {可放行的 late-load 命中\\异步镜像在 C0 读出原始字};
\node[b,fill=cyan!8] (reg) at (-2.0,0.0) {跨 EXU--LSU 寄存边界\\C1 扩展后供受限消费者};
\draw[a] (-2.0,3.55)--(mirror); \draw[a] (mirror)--(reg);

\node[b,fill=purple!7] (store) at (3.2,2.0) {store：主 BRAM 与镜像同址更新};
\node[b,fill=purple!7] (storepolicy) at (3.2,0.0) {整字写入并分配 tag\\窄写入则保守失效该行};
\draw[a] (3.2,3.55)--(store); \draw[a] (store)--(storepolicy);

\node[b,fill=gray!10] (miss) at (8.4,2.0) {load miss：访问下级数据存储};
\node[b,fill=gray!10] (response) at (8.4,0.0) {响应先完成本次 load};
\node[d] (epoch) at (8.4,-2.0) {store epoch 仍一致？};
\node[b,fill=gray!10,text width=4.15cm] (fill) at (5.7,-4.25) {是：回填主 BRAM、镜像与 tag};
\node[b,fill=gray!10,text width=4.15cm] (drop) at (11.1,-4.25) {否：仅交付响应，取消旧 miss 回填};
\draw[a] (8.4,3.55)--(miss); \draw[a] (miss)--(response); \draw[a] (response)--(epoch);
\draw[a] (epoch)--node[above left]{是}(fill); \draw[a] (epoch)--node[above right]{否}(drop);
\end{tikzpicture}}
\caption{D-cache 在命中、late-load、store 与 miss 情况下的处理流程}\end{figure}
